\section{Conclusiones}

Se logró implementar un sistema de chat multicliente basado en la arquitectura cliente-servidor, cubriendo
por completo el protocolo establecido en clase, además de incorporar una interfaz de usuario por terminal.

El servidor se desempeñó de manera bastante decente incluso bajo una carga considerable. Aunque esperaba
que el sistema colapsara con un número mucho menor de usuarios concurrentes, los resultados muestran que soportó
3,000 usuarios simultáneos con una tasa de error prácticamente nula, al menos en el escenario simulado.

Sin embargo, el proceso de desarrollo estuvo lleno de altibajos y de una carga de nuevos conocimientos
abismal. En retrospectiva, creo haber tomado algunas decisiones técnicas sobreestimando mi capacidad de
aprendizaje y el tiempo disponible. Por ejemplo, desarrollar el proyecto en dos lenguajes con los que 
tenía muy poca experiencia previa, o dedicar una cantidad considerable de tiempo a módulos no 
estrictamente necesarios, como el diseño detallado de la interfaz de terminal. 

Precisamente, el tener que aprender tantos conceptos y tecnologías nuevas a la vez limitó el tiempo 
disponible para el modelado inicial, lo cual me llevó a tomar decisiones subóptimas para resolver ciertos
problemas, mismas que terminé pagando con refactorizaciones.

Ahora que comprendo mejor las tecnologías con las que trabajé, espero poder modelar y anticipar mejor los próximos 
proyectos. En particular, pretendo delimitar el alcance desde el inicio y administrar mejor mis módulos de
trabajo, para no invertir tanto tiempo decidiendo en qué trabajar cada día.